\documentclass{fiche}
\metadonnees{3}{Voir et observer}{Bloc 1 — Le récit : présent et passé}

\begin{document}
\entete

\objectifs{%
\textbf{Langue} : present perfect et prétérit ; \emph{for}, \emph{since}, \emph{ago} ;
familles de mots.\par
\textbf{Texte} : Arthur Conan Doyle, \og A Scandal in Bohemia \fg (1891).\par
\textbf{Durée} : 1\,h\,30 — exercices 35 min, texte et questions 30 min, version 25 min.}

%% ===================================================================
\exercice[13 min]{Present perfect ou prétérit}

\rappel{\textbf{Rappel.} Le present perfect fait le lien avec maintenant : bilan, résultat,
expérience. Le prétérit coupe le lien : il exige un repère passé, exprimé ou sous-entendu.}

\consigne{Mettez le verbe entre parenthèses au present perfect ou au prétérit.}

\begin{anglais}
\begin{enumerate}[label=\arabic*.,itemsep=2.2mm,leftmargin=*]
  \item I \emph{(not see)} \trou{have not seen} Holmes since my marriage.
  \item He \emph{(solve)} \trou{solved} the Trepoff murder in 1887.
  \item ``You \emph{(put on)} \trou{have put on} seven and a half pounds since I saw you.''
  \item Watson \emph{(return)} \trou{returned} to civil practice a few months ago.
  \item ``How long \emph{(you / live)} \trou{have you lived} in Baker Street?'' --- ``For
    eight years.''
  \item I \emph{(pass)} \trou{passed} the well-remembered door last night.
  \item \emph{(you / ever / meet)} \trou{Have you ever met} a more perfect reasoning machine?
  \item He \emph{(not write)} \trou{has not written} to me since March.
  \item When \emph{(he / take)} \trou{did he take} the case?
  \item Holmes \emph{(just / light)} \trou{has just lit} a cigarette; he still \emph{(not
    say)} \trou{has not said} a word.
  \item The official police \emph{(abandon)} \trou{abandoned} those mysteries two years ago.
  \item Watson \emph{(stare)} \trou{stared} at the door and \emph{(wonder)} \trou{wondered}
    whether to ring.
\end{enumerate}
\end{anglais}

\note{Present perfect : 1, 3, 5, 7, 8, 10. Prétérit : 2, 4, 6, 9, 11, 12.}

\note[Le repère décide]{\emph{in 1887}, \emph{a few months ago}, \emph{last night},
\emph{two years ago}, \emph{when\ldots?} imposent le prétérit : le moment est daté.
\emph{since}, \emph{for}, \emph{ever}, \emph{just}, \emph{still not} appellent le present
perfect : le fait compte encore aujourd'hui. Phrase 3 : \emph{have put on} pour le bilan,
mais \emph{since I saw you} au prétérit, parce que la rencontre, elle, est datée.}

\note[Faute classique]{\emph{When have you arrived?} est impossible : \emph{when} demande un
moment précis, donc le prétérit. De même \emph{I have seen him yesterday}.}

%% ===================================================================
\exercice[10 min]{\emph{For}, \emph{since}, \emph{ago}}

\rappel{\textbf{Rappel.} \emph{For} + durée, \emph{since} + point de départ, \emph{ago} +
durée mais toujours avec le prétérit.}

\consigne{Complétez.}

\begin{anglais}
\begin{enumerate}[label=\arabic*.,itemsep=2.2mm,leftmargin=*]
  \item I have not seen my friend \trou{for} three years.
  \item He has lived in Baker Street \trou{since} 1881.
  \item She left London three years \trou{ago}.
  \item We have been friends \trou{since} our first case.
  \item Holmes has been at work \trou{for} several hours.
  \item The servant girl gave up her place a month \trou{ago}.
  \item Nobody has heard of him \trou{since} the spring.
  \item How long have you been scared of the dark? --- \trou{For} as long as I can remember.
  \item He has not touched the violin \trou{since} he took the case.
  \item They finished the investigation two days \trou{ago}, and the house has been still
    \trou{since} then.
\end{enumerate}
\end{anglais}

\note{\emph{Since} peut introduire un groupe nominal (\emph{since 1881}, \emph{since the
spring}) ou une proposition entière au prétérit (\emph{since he took the case}) : le point
de départ est alors raconté. \emph{Ago} se compte à rebours depuis maintenant et n'apparaît
jamais avec le present perfect.}

%% ===================================================================
\exercice[12 min]{Familles de mots}
\consigne{a) Complétez le tableau.}

\begin{center}\small
\begin{tabular}{@{}lll@{}}
\textsf{verbe} & \textsf{nom} & \textsf{adjectif} \\[1.5mm]
\emph{to deduce} & \trou{deduction} & \trou{deductive} \\[1mm]
\emph{to observe} & \trou{observation} & \trou{observant} \\[1mm]
\emph{to reason} & \trou{reasoning} & \trou{reasonable} \\[1mm]
\emph{to conclude} & \trou{conclusion} & \trou{conclusive} \\[1mm]
\emph{to suspect} & \trou{suspicion} & \trou{suspicious} \\[1mm]
\emph{to prove} & \trou{proof} & \trou{provable} \\[1mm]
\emph{to explain} & \trou{explanation} & \trou{explanatory} \\[1mm]
\emph{to observe} & \trou{observer} & --- \\
\end{tabular}
\end{center}

\note[Trois suffixes]{\emph{-tion} et \emph{-sion} forment des noms d'action à partir du
verbe (\emph{deduce $\rightarrow$ deduction}, \emph{conclude $\rightarrow$ conclusion}) ;
\emph{-ive} et \emph{-ous} des adjectifs (\emph{deductive}, \emph{suspicious}) ; \emph{-er}
le nom de celui qui fait l'action (\emph{observer}, \emph{reasoner}). Deux irrégularités à
retenir : \emph{to prove} donne \emph{proof} (avec \emph{f}), \emph{to explain} donne
\emph{explanation} (sans \emph{i}).}

\consigne{b) Complétez avec un mot du tableau.}
\begin{anglais}
\begin{enumerate}[label=\arabic*.,itemsep=2mm,leftmargin=*]
  \item Holmes is the most \trou{observant} man in London.
  \item His \trou{deduction} was simple once he had explained it.
  \item There is no \trou{proof} that she ever left the country.
  \item Watson found the \trou{explanation} ridiculously simple.
  \item Her behaviour that night was extremely \trou{suspicious}.
\end{enumerate}
\end{anglais}

%% ===================================================================
\newpage
\section{Texte}

\noindent\small\textit{Le docteur Watson, marié, a quitté l'appartement de Baker Street qu'il
partageait avec Sherlock Holmes.}\normalsize

\begin{texte}
I had seen little of Holmes lately. My marriage had drifted us away from each other. My own
complete happiness, and the home-centred interests which rise up around the man who first
finds himself master of his own establishment, were sufficient to absorb all my attention,
while Holmes, who \gl[to loathe]{loathed}{détester, avoir en horreur} every form of society
with his whole Bohemian soul, remained in our lodgings in Baker Street, buried among his old
books, and alternating from week to week between cocaine and ambition, the drowsiness of the
drug, and the fierce energy of his own keen nature. He was still, as ever, deeply attracted
by the study of crime, and occupied his immense faculties and extraordinary powers of
observation in following out those \gl[a clue]{clues}{un indice}, and clearing up those
mysteries which had been abandoned as hopeless by the official police. From time to time I
heard some vague account of his doings: of his \gl[a summons]{summons}{une convocation} to
Odessa in the case of the Trepoff murder, of his clearing up of the singular tragedy of the
Atkinson brothers at Trincomalee, and finally of the mission which he had accomplished so
delicately and successfully for the reigning family of Holland. Beyond these signs of his
activity, however, which I merely shared with all the readers of the daily press, I knew
little of my former friend and companion.

One night --- it was on the twentieth of March, 1888 --- I was returning from a journey to a
patient (for I had now returned to civil practice), when my way led me through Baker Street.
As I passed the well-remembered door, which must always be associated in my mind with my
\gl[a wooing]{wooing}{une cour amoureuse}, and with the dark incidents of the Study in
Scarlet, I was seized with a keen desire to see Holmes again, and to know how he was
employing his extraordinary powers. His rooms were brilliantly lit, and, even as I looked up,
I saw his tall, \gl[spare]{spare}{maigre, sec} figure pass twice in a dark silhouette against
the blind. He was pacing the room swiftly, eagerly, with his head sunk upon his chest and his
hands clasped behind him. To me, who knew his every mood and habit, his attitude and manner
told their own story. He was at work again. He had risen out of his drug-created dreams and
was hot upon the scent of some new problem. I rang the bell and was shown up to the chamber
which had formerly been in part my own.

His manner was not \gl[effusive]{effusive}{expansif}. It seldom was; but he was glad, I
think, to see me. With hardly a word spoken, but with a kindly eye, he waved me to an
armchair, threw across his case of cigars, and indicated a spirit case and a
\gl[a gasogene]{gasogene}{un siphon à eau de Seltz} in the corner. Then he stood before the
fire and looked me over in his singular introspective fashion.

``\gl[wedlock]{Wedlock}{le mariage} suits you,'' he remarked. ``I think, Watson, that you
have put on seven and a half pounds since I saw you.''

``Seven!'' I answered.

``Indeed, I should have thought a little more. Just a trifle more, I fancy, Watson. And in
practice again, I observe. You did not tell me that you intended to
\gl[to go into harness]{go into harness}{reprendre le collier}.''

``Then, how do you know?''

``I see it, I deduce it. How do I know that you have been getting yourself very wet lately,
and that you have a most clumsy and careless servant girl?''

``My dear Holmes,'' said I, ``this is too much. You would certainly have been burned, had you
lived a few centuries ago. It is true that I had a country walk on Thursday and came home in
a dreadful mess, but as I have changed my clothes I can't imagine how you deduce it. As to
Mary Jane, she is incorrigible, and my wife has given her notice, but there, again, I fail to
see how you work it out.''

He chuckled to himself and rubbed his long, nervous hands together.

``It is simplicity itself,'' said he; ``my eyes tell me that on the inside of your left shoe,
just where the firelight strikes it, the leather is \gl[to score]{scored}{entailler,
rayer} by six almost parallel cuts. Obviously they have been caused by someone who has very
carelessly scraped round the edges of the sole in order to remove crusted mud from it. Hence,
you see, my double deduction that you had been out in vile weather, and that you had a
particularly malignant boot-slitting specimen of the London
\gl[a slavey]{slavey}{une bonne à tout faire}. As to your practice, if a gentleman walks
into my rooms smelling of \gl[iodoform]{iodoform}{l'iodoforme, antiseptique}, with a black
mark of nitrate of silver upon his right forefinger, and a bulge on the right side of his
top-hat to show where he has \gl[to secrete]{secreted}{dissimuler} his stethoscope, I must be
dull, indeed, if I do not pronounce him to be an active member of the medical profession.''
\end{texte}
\source{Arthur Conan Doyle, \og A Scandal in Bohemia \fg, \emph{The Adventures of Sherlock
Holmes}, Londres, George Newnes, 1892.}

\partiel[15 min]{Questions}
\consigne{\en{Answer in English, in full sentences.}}

\begin{enumerate}[label=\arabic*.,itemsep=2mm,leftmargin=*]
  \item \en{Why have the two men seen so little of each other?}
    \lignes{2}
    \corr{Watson's marriage has drawn him away: his own happiness and his new home absorb
    him, while Holmes, who hates society, has stayed in Baker Street among his books.}
  \item \en{How does Watson know, from the street, that Holmes is working again?}
    \lignes{2}
    \corr{The rooms are brilliantly lit and he sees Holmes's tall thin silhouette pacing
    swiftly across the blind, head sunk on his chest and hands clasped behind him.}
  \item \en{What are the three things Holmes deduces about Watson, and from what clue?}
    \lignes{4}
    \corr{That he has put on weight (he simply looks him over); that he has been out in bad
    weather and has a careless servant, from the six parallel cuts on the inside of his left
    shoe, made by someone scraping off dried mud; that he is in practice again, from the
    smell of iodoform, the nitrate of silver on his forefinger and the bulge in his top-hat
    where the stethoscope is hidden.}
  \item \en{Watson says Holmes would have been burned a few centuries ago. What does he
    mean?}
    \lignes{2}
    \corr{That his deductions look like magic, and that in an earlier age he would have been
    taken for a sorcerer and burned as a witch.}
  \item \en{What is the difference between seeing and observing, according to this
    passage?}
    \lignes{3}
    \corr{Seeing is passive: Watson has looked at his own shoe and his own hat many times
    without drawing anything from them. Observing is reading a detail as a sign, and linking
    it to a cause.}
  \item \en{Pick out two details that show Watson admires Holmes and one that shows he
    disapproves of him.}
    \lignes{3}
    \corr{Admiration : \emph{the most perfect reasoning and observing machine},
    \emph{his immense faculties and extraordinary powers of observation}, \emph{I could not
    help laughing at the ease with which he explained}. Réserve : \emph{alternating\ldots
    between cocaine and ambition}, \emph{his drug-created dreams}.}
\end{enumerate}

\note[Temps du récit]{Watson raconte au prétérit (\emph{I rang the bell}, \emph{he stood
before the fire}) et remonte plus haut au pluperfect (\emph{I had seen little of Holmes},
\emph{my marriage had drifted us away}, \emph{which had been abandoned}). Mais dans les
répliques, au discours direct, Holmes parle de maintenant : \emph{you have put on},
\emph{you have been getting yourself wet}, \emph{my wife has given her notice}. Le partage
entre les deux temps suit exactement le partage entre récit et parole.}

%% ===================================================================
\newpage
\section{Version}
\consigne{Traduisez le passage en français. C'est le premier paragraphe de la nouvelle : il
précède immédiatement le texte que vous venez de lire.}

\begin{version}
To Sherlock Holmes she is always \emph{the} woman. I have seldom heard him mention her under
any other name. In his eyes she eclipses and predominates the whole of her sex. It was not
that he felt any emotion akin to love for Irene Adler. All emotions, and that one
particularly, were \gl[abhorrent]{abhorrent}{odieux, insupportable} to his cold, precise but
admirably balanced mind. He was, I take it, the most perfect reasoning and observing machine
that the world has seen, but as a lover he would have placed himself in a false position. He
never spoke of the softer passions, save with a \gl[a gibe]{gibe}{une raillerie} and a
\gl[a sneer]{sneer}{un ricanement méprisant}. They were admirable things for the observer ---
excellent for drawing the veil from men's motives and actions. But for the trained reasoner
to admit such intrusions into his own delicate and finely adjusted temperament was to
introduce a distracting factor which might throw a doubt upon all his mental results.
\gl[grit]{Grit}{du gravier, des impuretés} in a sensitive instrument, or a crack in one of his
own high-power lenses, would not be more disturbing than a strong emotion in a nature such as
his. And yet there was but one woman to him, and that woman was the late Irene Adler, of
dubious and questionable memory.
\end{version}
\source{\emph{Ibid.}}

\lignes{22}

\corr{\textbf{Proposition de traduction.} Pour Sherlock Holmes, elle est toujours \emph{la}
femme. Je l'ai rarement entendu la désigner autrement. À ses yeux, elle éclipse et domine
tout son sexe. Non qu'il eût éprouvé pour Irene Adler quelque sentiment voisin de l'amour.
Toutes les émotions, et celle-là en particulier, répugnaient à son esprit froid, précis,
admirablement équilibré. Il était, je crois, la plus parfaite machine à raisonner et à
observer que le monde ait connue ; mais en amoureux, il se fût mis dans une situation fausse.
Jamais il ne parlait des passions tendres, sinon par une raillerie ou un ricanement. C'étaient
là d'admirables objets pour l'observateur, excellents pour lever le voile sur les mobiles et
les actes des hommes ; mais pour le raisonneur exercé, admettre de telles intrusions dans son
tempérament délicat et finement réglé, c'était y introduire un facteur de trouble propre à
jeter le doute sur tous ses résultats. Du gravier dans un instrument de précision, une fêlure
dans l'une de ses puissantes lentilles, ne le dérangeraient pas davantage qu'une émotion forte
dans une nature comme la sienne. Et pourtant il n'y avait pour lui qu'une seule femme, et
cette femme était feu Irene Adler, de douteuse et discutable mémoire.}

\note[Choix de traduction 1]{Le paragraphe commence au présent (\emph{she is always the
woman}, \emph{I have seldom heard}) puis bascule au prétérit (\emph{It was not that he
felt}) : Watson écrit après coup, et le présent est celui de son jugement d'aujourd'hui. Il
faut garder les deux plans, sans tout mettre au passé.}

\note[Choix de traduction 2]{\emph{It was not that he felt any emotion akin to love} : le
tour \emph{it is not that\ldots} se rend par \og non que \fg{} + subjonctif. \emph{Any} dans
une phrase négative n'est pas \og n'importe quel \fg{} mais \og le moindre \fg.}

\note[Choix de traduction 3]{\emph{as a lover he would have placed himself in a false
position} : \emph{as} = \og en tant que \fg, \og en \fg. \emph{Would have} + participe est
un irréel du passé : \og il se serait mis \fg{} ou, en langue soutenue, \og il se fût mis \fg.}

\note[Choix de traduction 4]{\emph{the most perfect\ldots machine that the world has seen} :
après un superlatif, l'anglais emploie le present perfect, le français le subjonctif
(\og que le monde ait connue \fg). C'est une correspondance à mémoriser.}

\note[Choix de traduction 5]{\emph{save with a gibe and a sneer} : \emph{save} est ici une
préposition littéraire, \og sinon \fg, \og excepté \fg, et non le verbe \og sauver \fg. De
même \emph{but one woman} : \emph{but} vaut \og seulement \fg.}

\note[Choix de traduction 6]{\emph{Grit in a sensitive instrument\ldots would not be more
disturbing than} : la comparaison porte sur un sujet très long ; le français supporte mal
d'attendre le verbe aussi longtemps. On peut découper, ou garder la construction en allégeant
les groupes, comme ci-dessus.}

%% ===================================================================
\partiel{Lexique à mémoriser}
\begin{itemize}[label=--,itemsep=0.8mm,leftmargin=*]\small
  \item \textbf{a clue} — un indice
  \item \textbf{to deduce} — déduire ; \textit{nom} a deduction
  \item \textbf{to observe} — observer, remarquer ; \textit{noms} observation, an observer
  \item \textbf{proof} — la preuve ; \textit{verbe} to prove\textit{, avec} v/f \textit{comme} belief / believe
  \item \textbf{keen} — vif, aiguisé ; keen on \textit{« passionné de »}
  \item \textbf{eager} — impatient, avide
  \item \textbf{to loathe} — détester ; \textit{plus fort que} to dislike
  \item \textbf{seldom} — rarement ; \textit{terme littéraire pour} rarely
  \item \textbf{clumsy} — maladroit
  \item \textbf{careless} — négligent ; care + -less\textit{, contraire de} careful
  \item \textbf{a mess} — un désordre, un gâchis
  \item \textbf{to pace} — arpenter
  \item \textbf{to chuckle} — glousser, rire sous cape
  \item \textbf{vile} — infect, exécrable
  \item \textbf{mud} — la boue ; \textit{adjectif} muddy
  \item \textbf{a sole} — une semelle
  \item \textbf{to fail to} — ne pas parvenir à
  \item \textbf{akin to} — voisin de, apparenté à ; \textit{de} kin \textit{« la parenté »}
  \item \textbf{a veil} — un voile ; to draw the veil from \textit{« lever le voile sur »}
  \item \textbf{the late} — feu, défunt ; the late Irene Adler
\end{itemize}

\end{document}
